% =====================================================================
% What Actually Limits Bangla Image Captioning?
% A Controlled LoRA Study with an Honest Evaluation Framework
%
% Draft manuscript (working copy, 2026-07-18). Compiles with pdfLaTeX
% (e.g., Overleaf default). Bangla examples appear in romanized form
% with English glosses; the submission version will typeset Bangla
% script via XeLaTeX + a Bengali font.
%
% Human-evaluation pilot results are PENDING and marked as such.
% =====================================================================
\documentclass[11pt]{article}

\usepackage[margin=1in]{geometry}
\usepackage{booktabs}
\usepackage{amsmath}
\usepackage{graphicx}
\usepackage{multirow}
\usepackage{url}
\usepackage[hidelinks]{hyperref}
\usepackage{xcolor}

\newcommand{\pending}[1]{\textcolor{red!70!black}{[PENDING: #1]}}
\newcommand{\semmargin}{\textsc{SemMargin}}

\title{What Actually Limits Bangla Image Captioning?\\
\large A Controlled LoRA Study with an Honest Evaluation Framework}

\author{Onik Jahan Sagor\\
\small \textit{Draft manuscript --- July 2026}}

\date{}

\begin{document}
\maketitle

\begin{abstract}
Bangla, spoken by roughly 270 million people, remains poorly served by
modern vision--language models (VLMs), and published Bangla captioning
results are difficult to interpret: they span an order of magnitude on
the same corpora, rarely name their evaluation protocol, and have never
been checked against trivial baselines. We present a controlled study
that separates the real bottlenecks from the measurement noise. First,
through a six-stage ablation ladder on GiT-base, we show that
\emph{vocabulary coverage}, not trainable capacity, decides whether a
VLM can produce Bangla at all: attention-only LoRA, head retraining,
and full embedding retraining all fail identically until the tokenizer
is extended with 29{,}055 Bangla wordpieces, at which point Bangla
generation emerges immediately. The converse experiment on
Qwen2-VL-2B---whose vocabulary already covers Bangla but which produces
zero Bangla captions zero-shot---shows the same attention-only LoRA
recipe unlocking fluent Bangla within 120 training samples. Second, we
show that caption \emph{composition} is bound by decoder depth and
output \emph{diversity} by data scale, and that a two-stage curriculum
(machine-translated-then-post-edited silver corpus, followed by a short
native-gold finetune) yields our best system: BLEU-4 0.108 and CIDEr
0.354 on BAN-Cap with 0.178\% trainable parameters and under seven
hours on a consumer 8\,GB GPU. Third, we audit the evaluation practice
of the field: no-vision baselines match or beat trained models on
BLEU-1 and BERTScore; multilingual CLIPScore assigns its highest score
to output containing no Bangla whatsoever; and a faithful
reproduction of the strongest published BanglaLekha system (reported
BLEU-4 0.408) reaches only 0.064 under the paper's own stated protocol
--- while the scorer implementation alone moves BLEU-4 by $+68\%$. We
introduce \semmargin{}, a corpus-prior-corrected semantic similarity
metric that scores all no-vision baselines at exactly zero by
construction, and we release a blinded human-evaluation protocol for
Bangla captioning. All experiments are reproducible from configuration
files on a single consumer GPU.
\end{abstract}

% =====================================================================
\section{Introduction}
% =====================================================================

Image captioning for high-resource languages has been transformed by
large pre-trained vision--language models, but Bangla --- among the
world's most spoken languages --- still relies on small CNN--LSTM and
CNN--transformer pipelines trained from scratch on corpora of fewer
than ten thousand images. Reported results are hard to build on: BLEU-4
scores on the BanglaLekha corpus range from $2.2\times10^{-308}$
\cite{palash2021} to 0.408 \cite{bornon2021} for architecturally
similar systems, no published Bangla captioning paper names its BLEU
implementation, and none reports a no-vision baseline.

This paper asks two questions. \textbf{(Q1) What actually limits a
modern VLM's ability to caption images in Bangla?} We answer with a
chain of controlled experiments that isolate three walls ---
vocabulary coverage, decoder depth, and data scale --- and a
curriculum recipe that addresses the third on a consumer GPU.
\textbf{(Q2) Can the standard evaluation practice of the field detect
progress at all?} We answer mostly in the negative, with three
experiments --- trivial-baseline flooring, a published-score audit
with a full recipe reproduction, and a metric-blindness study --- and
we contribute the corrective: a specificity-corrected similarity
metric (\semmargin{}) and a blinded human-evaluation protocol.

Our contributions:
\begin{enumerate}
  \item \textbf{Mechanism (C1).} A six-stage ablation ladder
  establishing that vocabulary coverage is necessary and sufficient
  for Bangla generation capability in a frozen-backbone LoRA setting,
  demonstrated in both directions (GiT-base without coverage;
  Qwen2-VL-2B with it). A vocabulary/embedding bridge recipe (tokenizer
  extension + donor embedding initialization) that unlocks Bangla in a
  model with essentially no Bangla tokens.
  \item \textbf{Recipe (C2).} A QLoRA + curriculum recipe --- silver
  pretraining on a 4$\times$ machine-translated-then-post-edited
  corpus followed by a short gold finetune --- that improves every
  text metric over gold-only training with half the gold budget
  (BLEU-4 $0.082 \to 0.108$, CIDEr $0.300 \to 0.354$) while retaining
  output-vocabulary diversity, at a total cost of $\sim$7 GPU-hours on
  an 8\,GB RTX 4060 Ti.
  \item \textbf{Evaluation audit (C3).} The first trivial-baseline
  flooring of Bangla captioning metrics; a documented protocol audit
  of the published literature; and, to our knowledge, the first full
  reproduction attempt of a published Bangla captioning system, which
  fails to reproduce by a factor of 6.4$\times$ under the original
  paper's own stated protocol.
  \item \textbf{Metric + benchmark (C4).} \semmargin{}, a
  corpus-prior-corrected semantic similarity metric that zeroes all
  no-vision baselines by construction while reproducing the system
  ranking of the full n-gram study; and a blinded adequacy/fluency
  human-evaluation protocol with a public annotation tool.
  \pending{pilot human ratings and metric--human correlation table}
\end{enumerate}

% =====================================================================
\section{Related Work}
% =====================================================================

\paragraph{Bangla image captioning.} Chittron \cite{chittron2019}
introduced the task with a VGG16+LSTM pipeline and, notably, its
authors already flagged BLEU as ``admittedly extremely poor'' with a
single reference. Subsequent work --- Bornon \cite{bornon2021}
(InceptionV3 + transformer), CNN--Transformer \cite{palash2021}
(ResNet-101 + transformer), BAN-Cap's own adaptive-attention baseline
\cite{bancap2022}, and BiCap \cite{bicap2025} --- evaluates exclusively
with n-gram metrics on one- or two-reference corpora. Section
\ref{sec:audit} examines this literature quantitatively.

\paragraph{PEFT for multilingual VLMs.} PALO \cite{palo2024}, mBLIP
\cite{mblip2023}, Chitrarth \cite{chitrarth2025}, and Chitranuvad
\cite{chitranuvad2025} extend VLMs to Indic and other languages at
scale, generally with backbones whose tokenizers already cover the
target languages (or with full finetuning). None studies Bangla in a
dedicated PEFT setting, and none isolates vocabulary coverage as a
variable; our ablation ladder does exactly that. LoRA
\cite{lora2021} and QLoRA \cite{qlora2023} provide the adaptation
machinery.

\paragraph{Captioning evaluation.} Critiques of n-gram captioning
metrics are long-standing \cite{kilickaya2017}. Reference-free scoring
via CLIPScore \cite{clipscore2021} and learned metrics such as Polos
\cite{polos2024} are validated only for English. The
\emph{distinctiveness} line of work --- CIDErBtw \cite{ciderbtw2020}
and the group embedding gap \cite{geg2022} --- contrasts a caption's
similarity to its own image against similar images; \semmargin{}
adapts this device to the reference-text side, over the full
evaluation corpus, as a general quality metric for settings where
image--text encoders are unreliable (which, as we show for Bangla,
they are).

% =====================================================================
\section{The Vocabulary Wall}
\label{sec:vocab}
% =====================================================================

\subsection{Tokenizer audit}

We measured tokenizer fertility (subword pieces per word) on 2{,}000
Bangla Wikipedia sentences across nine tokenizers, and re-ran the audit
on BAN-Cap captions. GiT-base \cite{git2022} inherits
\texttt{bert-base-uncased}: only 72 of its 30{,}522 tokens contain any
Bangla codepoint, fertility is 4.88, and 2.5\% of tokens are unknown.
BanglaBERT \cite{banglabert2022} scores 1.35 on the same text. Our
bridged tokenizer (below) reaches 1.33 on Wikipedia and exactly ties
BanglaBERT (1.176) on BAN-Cap captions. Byte-level BPE tokenizers
(Qwen2-VL, Florence-2) report zero Bangla vocabulary by codepoint
counting --- a metric artifact of byte encoding; for those, fertility
and round-trip fidelity are the meaningful measures.

\subsection{A six-stage ablation ladder}

All stages use GiT-base with identical data and decoding; only the
adaptation surface changes.

\begin{enumerate}
  \item \textbf{Full finetuning} converges (loss $12.3 \to 5.1$) but
  generates only a malformed repetition of the single Bangla-adjacent
  wordpiece in the vocabulary.
  \item \textbf{Attention-only LoRA} (0.125\% trainable): the model
  answers in \emph{English} --- with embeddings and head frozen it
  cannot even reach its few Bangla tokens.
  \item \textbf{LoRA + trained LM head}: loss drops to 3.46; output
  identical to full finetuning's malformed repetition.
  \item \textbf{LoRA + head + embeddings} (+24M trainable parameters):
  output \emph{identical} to stage 3. No weight adjustment can create
  tokens that do not exist.
  \item \textbf{The bridge}: we extend the vocabulary with 29{,}055
  BanglaBERT wordpieces ($30{,}522 \to 59{,}577$), initializing new
  input embeddings from the donor model. Generation immediately flips
  to real (if trivial) Bangla.
  \item \textbf{The converse, on a modern VLM
  (Sec.~\ref{sec:qwen})}: Qwen2-VL-2B's byte-BPE vocabulary covers
  Bangla, yet zero-shot it produces \textbf{0 of 809} Bangla captions.
  The \emph{same attention-only LoRA that failed in stage 2} ---
  0.178\% trainable, vision tower and embeddings frozen --- produces
  fluent Bangla within $\sim$120 training samples and 809/809 Bangla
  outputs after training.
\end{enumerate}

Together, stages 1--6 close the argument in both directions:
\textbf{vocabulary coverage is the deciding factor} between the two
failure modes, weight capacity is not the bottleneck, and the bridge
buys coverage when the base model lacks it.

% =====================================================================
\section{The Depth and Data Walls}
\label{sec:depth}
% =====================================================================

\subsection{Composition is decoder-bound}

The bridged GiT trained on real corpora produces genuine but
\emph{templated} Bangla. Three post-hoc diagnostics on the trained
adapters locate the cause. (i) \emph{Image scrambling}: replacing
inputs with gray/noise/shuffled images shows the model takes nouns
from the image but not sentence structure. (ii) \emph{Forced-length
decoding} triples output length without producing grammar, refuting a
``hidden composition'' hypothesis. (iii) \emph{Per-reference
cross-entropy} exposed an evaluation bias in the two-reference
BanglaLekha corpus (3.76$\times$ perplexity gap between short and long
references) that our corpus migration to BAN-Cap \cite{bancap2022}
(five native references per image) structurally removed
(1.14$\times$).

Budget and placement controls confirm the diagnosis: doubling training
steps and adding image-tower LoRA each improve unigram precision
(BLEU-1 up to 0.498) but leave BLEU-4 at 0.000. Only the 28-layer
decoder of Qwen2-VL-2B breaks the composition wall (BLEU-4
$0.000 \to 0.082$), at 0.178\% trainable parameters.

\subsection{Diversity is data-bound, and a curriculum fixes it}
\label{sec:qwen}

The BAN-Cap-trained Qwen adapter collapses to \textbf{494 unique
output tokens} across the 809-image validation set (references use
4{,}797); top-p sampling does not recover the missing vocabulary, so
the collapse is data-limited, not decode-limited. We therefore
pretrained the same recipe on BanglaView --- 31{,}783 Flickr30k images
with five machine-translated-then-native-post-edited captions each,
4$\times$ the image scale --- and then resumed the adapter for 1{,}500
steps on native BAN-Cap.

\begin{table}[t]
\centering
\caption{BAN-Cap validation (809 images, 5 native references).
Curriculum = BanglaView silver pretrain (6k steps) + gold finetune
(1.5k steps). Unique tokens counts distinct words across all 809
predictions (references: 4{,}797).}
\label{tab:curriculum}
\small
\begin{tabular}{lccccccc}
\toprule
Adapter & B1 & B2 & B3 & B4 & CIDEr & BERTScore-F1 & Uniq.\ tokens \\
\midrule
BAN-Cap-only (3k steps) & 0.532 & 0.331 & 0.177 & 0.082 & 0.300 & 0.804 & 494 \\
BanglaView-only (6k steps) & 0.483 & 0.292 & 0.162 & 0.080 & 0.264 & 0.780 & \textbf{699} \\
\textbf{Curriculum} & \textbf{0.558} & \textbf{0.355} & \textbf{0.203} & \textbf{0.108} & \textbf{0.354} & \textbf{0.812} & 554 \\
\midrule
Published BAN-Cap baseline \cite{bancap2022} & 0.738 & 0.495 & 0.329 & 0.208 & 0.308 & --- & --- \\
\bottomrule
\end{tabular}
\end{table}

Table~\ref{tab:curriculum}: the curriculum wins \emph{every} text
metric over gold-only training --- BLEU-4 $+32\%$, CIDEr $+18\%$ ---
with half the gold budget, and retains a $+12\%$ vocabulary advantage
over the gold-only collapse. Cross-corpus transfer is nearly free
(BanglaView-only matches in-domain BLEU-4 on a corpus it never saw).
The published BAN-Cap row is a full-finetuned, task-specific
architecture; our system reaches roughly half its BLEU-4 with 0.178\%
trainable parameters and $\sim$7 hours on an 8\,GB consumer GPU ---
and Section~\ref{sec:audit} gives strong reasons to treat all
cross-paper comparisons with caution.

% =====================================================================
\section{Can the Field's Evaluation Detect Progress?}
\label{sec:audit}
% =====================================================================

\subsection{No-vision baselines floor the metrics}

We scored four trivial baselines that never see the image ---
\texttt{constant} (most frequent training caption), \texttt{mode\_len}
(most frequent caption of median length), \texttt{random} (random
training caption), \texttt{freq\_words} (the eight most frequent
training tokens strung together) --- under the identical multi-reference
battery as every real system.

\begin{table}[t]
\centering
\caption{Selected rows, BAN-Cap val (809 $\times$ 5 refs). A bag of
frequent words beats real vision systems on BLEU-1 and BERTScore; only
CIDEr and \semmargin{} (Sec.~\ref{sec:semmargin}) resist.}
\label{tab:trivial}
\small
\begin{tabular}{llcccccc}
\toprule
System & Sees image? & B1 & B4 & CIDEr & BERTScore-F1 & M-CLIPScore \\
\midrule
\texttt{freq\_words} & no & 0.457 & 0.000 & 0.104 & 0.760 & 0.542 \\
\texttt{mode\_len} & no & 0.198 & 0.020 & 0.042 & 0.752 & 0.539 \\
GiT+bridge (best) & yes & 0.498 & 0.000 & 0.172 & 0.758 & 0.546 \\
Qwen2-VL zero-shot (0\% Bangla) & yes & 0.000 & 0.000 & 0.000 & 0.690 & \textbf{0.803} \\
Qwen QLoRA curriculum & yes & \textbf{0.558} & \textbf{0.108} & \textbf{0.354} & \textbf{0.812} & 0.560 \\
\bottomrule
\end{tabular}
\end{table}

Three metric indictments follow from Table~\ref{tab:trivial}.
\textbf{BLEU-1 is frequency-blind}: a fixed word salad scores 0.457.
\textbf{BERTScore is template-blind}: on the two-reference BanglaLekha
protocol a constant three-word caption reaches BERTScore-F1
\textbf{0.798} --- within 0.006 of our best real system's BAN-Cap
score. \textbf{M-CLIPScore is language-blind}: it assigns its highest
score (0.803) to the zero-shot system whose 809 outputs contain no
Bangla codepoint at all. Only CIDEr's tf-idf weighting resists
common-word gaming, and even it leaks 0.104 to the no-vision band. Our
own GiT rows do \emph{not} clear the no-vision band on n-gram metrics
--- an honest negative we believe every captioning paper should be
required to report.

\subsection{The published literature does not reproduce}

A full-text audit of published Bangla captioning results found:
sentence-level and corpus-level BLEU mixed freely; a published BLEU-4
of $2.2\times10^{-308}$ (an unsmoothed zero) in a paper whose abstract
claims BLEU-2 of 0.630 \cite{palash2021}; reference-count effects worth
$0.738 \to 0.480$ BLEU-1 on the same model (BAN-Cap's own ablation
\cite{bancap2022}); vocabulary truncation to the top 5{,}000
\cite{bornon2021} or top 1{,}500 words; and cross-paper baseline rows
that do not match the cited papers.

We went one step further and \textbf{reproduced the strongest
published BanglaLekha system}: Bornon's transformer \cite{bornon2021}
(reported B1/B2/B3/B4 = 0.665/0.556/0.476/\textbf{0.408}). No code or
weights were released; we reimplemented the paper's stated recipe
exactly --- InceptionV3 features, 3-layer/1-head transformer, top-5{,}000
Keras tokenizer, 7{,}154/1{,}000/1{,}000 split, 50 epochs, Adam with
the Vaswani schedule --- trained it to full convergence, and scored the
same predictions three ways (Table~\ref{tab:bornon}).

\begin{table}[t]
\centering
\caption{Bornon reproduction on BanglaLekha (2 refs, 1{,}000 test
images). The same trained model, three scoring protocols.}
\label{tab:bornon}
\small
\begin{tabular}{lcccc}
\toprule
Scoring protocol & B1 & B2 & B3 & B4 \\
\midrule
Published \cite{bornon2021} & 0.665 & 0.556 & 0.476 & 0.408 \\
Ours --- their stated protocol (NLTK, top-5000) & 0.374 & 0.206 & 0.111 & 0.064 \\
Ours --- same scorer, no truncation & 0.367 & 0.203 & 0.109 & 0.063 \\
Ours --- pycocoevalcap & 0.452 & 0.288 & 0.172 & 0.106 \\
\bottomrule
\end{tabular}
\end{table}

Findings: (i) the published number does not reproduce from the
information the paper provides --- a 6.4$\times$ BLEU-4 gap that no
protocol variant we tested closes; (ii) the scorer implementation
alone (NLTK vs.\ pycocoevalcap) moves BLEU-4 by $+68\%$ on identical
predictions --- a free parameter in every published number, since no
paper names its implementation; (iii) the suspected top-5{,}000
truncation mechanism is \emph{innocent} on this corpus (its fitted
vocabulary is only 5{,}075 words) --- the audit kills one of its own
hypotheses; and (iv) under honest scoring the reproduced architecture
is actually the \emph{best} honest BanglaLekha system we measured
(CIDEr 0.495, 2.5$\times$ the trivial floor) --- the field's models
were better than its numbers were meaningful. We state explicitly:
this is a recipe reproduction under documented assumptions, not a
rerun of the original artifact, and it accuses no one; it measures
what the published information suffices to reproduce.

\subsection{\semmargin{}: subtracting the corpus prior}
\label{sec:semmargin}

Generic captions defeat embedding-similarity metrics because they sit
near the semantic centroid of the corpus. We correct for this
directly. With $e(\cdot)$ a sentence encoder (LaBSE
\cite{labse2022}), hypothesis $h_i$ for image $i$, and reference sets
$R_j$:
\begin{align}
\mathrm{raw}(h_i) &= \tfrac{1}{|R_i|}\textstyle\sum_{r \in R_i} \cos(e(h_i), e(r)),\\
\mathrm{prior}(h_i) &= \tfrac{1}{n-1}\textstyle\sum_{j \neq i} \tfrac{1}{|R_j|}\sum_{r \in R_j} \cos(e(h_i), e(r)),\\
\semmargin(h_i) &= \mathrm{raw}(h_i) - \mathrm{prior}(h_i).
\end{align}
A caption equally similar to every image's references --- any generic
caption --- scores zero \emph{by construction}. The device follows the
distinctiveness literature (CIDErBtw \cite{ciderbtw2020}, GEG
\cite{geg2022}); our formulation is reference-text-side, over the full
evaluation corpus, with no image encoder (unusable for Bangla, per the
M-CLIPScore result) and no curated similar-image sets.

\begin{table}[t]
\centering
\caption{\semmargin{} on BAN-Cap val. ``rank'' is the mean percentile
of the own-image similarity among all images (reference-side
self-retrieval). Human ceiling: each first reference scored against
the remaining four.}
\label{tab:semmargin}
\small
\begin{tabular}{lcccc}
\toprule
System & raw & prior & margin & rank \\
\midrule
Trivial baselines (all four) & 0.28--0.48 & = raw & \textbf{0.000} & 0.50 \\
GiT+bridge (three variants) & 0.53 & 0.38 & 0.145--0.153 & 0.91 \\
Qwen QLoRA BAN-Cap & 0.644 & 0.466 & 0.177 & 0.937 \\
Qwen QLoRA curriculum & 0.661 & 0.471 & \textbf{0.190} & 0.953 \\
Human (held-out reference) & 0.647 & 0.437 & 0.210 & 0.968 \\
Qwen2-VL zero-shot (0\% Bangla) & 0.601 & 0.361 & 0.240 & 0.980 \\
\bottomrule
\end{tabular}
\end{table}

Table~\ref{tab:semmargin}: the raw column shows plain embedding cosine
--- the uncorrected metric --- giving templates 0.43--0.48, nearly
indistinguishable from real systems; the margin column zeroes them
exactly. The margin reproduces the system ordering of the full n-gram
study with no shared machinery, and the curriculum system reaches 91\%
of the human held-out-reference margin. The one deliberate failure ---
the English-output zero-shot row tops the table because LaBSE is
cross-lingual by design --- demonstrates that \semmargin{} measures
the \emph{specificity} axis only, and must be paired with a trivially
computable language-fidelity gate. We argue this three-axis
decomposition (language fidelity, specificity, fluency/adequacy) is
the correct shape for low-resource captioning evaluation: every metric
failure documented above is a collapse of one axis into another.

\subsection{Human evaluation}

We built a blinded protocol: 20 seeded-random BAN-Cap validation
images $\times$ 4 systems (GiT image-tower LoRA, Qwen QLoRA gold-only,
Qwen curriculum, and a trivial template as attention check), captions
shuffled per image, system identities withheld, adequacy and fluency
on anchored 1--5 Likert scales, collected through a purpose-built web
annotation tool with reviewer-uniqueness checks. Ratings feed
per-system means, Krippendorff's $\alpha$, and per-caption
metric--human correlations for every metric in this paper.

\pending{Pilot ratings are being collected at the time of this draft.
The correlation table --- including whether \semmargin{} tracks human
adequacy better than BLEU/BERTScore, as its construction predicts ---
will complete Section 5. The full C3 benchmark (200 images $\times$ 3
annotators) is planned as the scaled follow-up.}

% =====================================================================
\section{Efficiency}
% =====================================================================

All results were produced on one RTX 4060 Ti (8\,GB). Representative
wall-clock: Qwen2-VL-2B QLoRA gold-only, 2.7\,h; BanglaView silver
stage, 5.0\,h; gold finetune stage, 78\,min; full 809-image scoring
battery, $\sim$35\,min; the Bornon reproduction trains in 13\,min. The
complete curriculum system costs $\sim$7 GPU-hours --- the point being
that honest, reproducible Bangla captioning research is feasible
without institutional compute.

% =====================================================================
\section{Limitations}
% =====================================================================

Our best BLEU-4 (0.108) remains below BAN-Cap's published full-FT
baseline (0.208) --- though Section~\ref{sec:audit} shows why
cross-paper numbers deserve scrutiny, ours included. The Bornon
reproduction rests on documented assumptions where the paper is
silent (d\_model, decoding, seed); its conclusion is about
reproducibility-from-published-information, not about the original
artifact. \semmargin{} requires a multilingual sentence encoder and
inherits its biases; its human validation is pending. The human-eval
pilot is small (20 images) and its annotator pool is a convenience
sample of native speakers. BanglaView's silver captions are
machine-translated then post-edited; translation artifacts may shape
the curriculum gains. GiT-base and BLIP-base share a tokenizer, so our
audit's coverage claims there are one measurement, not two.

% =====================================================================
\section{Conclusion}
% =====================================================================

Three walls limit Bangla image captioning --- vocabulary coverage,
decoder depth, and data scale --- and none of them is architecture.
Each has a cheap, reproducible mitigation: a tokenizer/embedding
bridge, a deeper frozen decoder under QLoRA, and a silver-then-gold
curriculum, jointly costing hours on a consumer GPU. The field could
not see these walls because its measurement practice cannot
distinguish a trained model from a frequency table: trivial baselines
floor the standard metrics, the strongest published number does not
reproduce, and the scorer itself is an unnamed free parameter. We
offer the corrective we could build --- trivial-baseline flooring as
mandatory practice, a specificity-corrected metric with a proven
zero-floor, and a blinded human protocol --- and release every
config, script, and result.

% =====================================================================
\begin{thebibliography}{99}
\small

\bibitem{chittron2019} M.~Rahman, N.~Mohammed, N.~Mansoor, S.~Momen.
Chittron: An automatic Bangla image captioning system.
\emph{Procedia Computer Science} 154 (2019).

\bibitem{bornon2021} F.~M.~Shah, M.~Humaira, M.~A.~R.~K.~Jim,
A.~S.~Ami, S.~Paul. Bornon: Bengali image captioning with
transformer-based deep learning approach. \emph{SN Computer Science}
(2021). arXiv:2109.05218.

\bibitem{palash2021} M.~A.~H.~Palash et al. Bangla image caption
generation through CNN-transformer based encoder-decoder network.
arXiv:2110.12442 (2021).

\bibitem{bancap2022} M.~F.~Khan et al. BAN-Cap: A multi-purpose
English-Bangla image descriptions dataset. \emph{LREC} (2022).
arXiv:2205.14462.

\bibitem{bicap2025} BiCap: Bilingual image captioning for Bangla.
\emph{BLP Workshop} (2025). aclanthology.org/2025.banglalp-1.6.

\bibitem{lora2021} E.~Hu et al. LoRA: Low-rank adaptation of large
language models. \emph{ICLR} (2022). arXiv:2106.09685.

\bibitem{qlora2023} T.~Dettmers et al. QLoRA: Efficient finetuning of
quantized LLMs. \emph{NeurIPS} (2023). arXiv:2305.14314.

\bibitem{git2022} J.~Wang et al. GiT: A generative image-to-text
transformer for vision and language. \emph{TMLR} (2022).
arXiv:2205.14100.

\bibitem{qwen2vl2024} P.~Wang et al. Qwen2-VL: Enhancing
vision-language model's perception of the world at any resolution.
arXiv:2409.12191 (2024).

\bibitem{banglabert2022} A.~Bhattacharjee et al. BanglaBERT: Language
model pretraining and benchmarks for low-resource language
understanding evaluation in Bangla. \emph{Findings of NAACL} (2022).

\bibitem{palo2024} M.~Maaz et al. PALO: A polyglot large multimodal
model for 5B people. arXiv:2402.14818 (2024).

\bibitem{mblip2023} G.~Geigle et al. mBLIP: Efficient bootstrapping of
multilingual vision-LLMs. arXiv:2307.06930 (2023).

\bibitem{chitrarth2025} Chitrarth: Bridging vision and language for a
billion people. arXiv:2502.15392 (2025).

\bibitem{chitranuvad2025} Chitranuvad: Adapting multi-lingual LLMs for
multimodal translation. arXiv:2502.20420 (2025).

\bibitem{kilickaya2017} M.~Kilickaya et al. Re-evaluating automatic
metrics for image captioning. \emph{EACL} (2017).

\bibitem{clipscore2021} J.~Hessel et al. CLIPScore: A reference-free
evaluation metric for image captioning. \emph{EMNLP} (2021).

\bibitem{polos2024} Y.~Wada et al. Polos: Multimodal metric learning
from human feedback for image captioning. \emph{CVPR} (2024).

\bibitem{ciderbtw2020} J.~Wang, W.~Xu, Q.~Wang, A.~B.~Chan. Compare
and reweight: Distinctive image captioning using similar images sets.
\emph{ECCV} (2020).

\bibitem{geg2022} Y.~Zhang et al. Distinctive image captioning via
CLIP guided group optimization. arXiv:2208.04254 (2022).

\bibitem{labse2022} F.~Feng et al. Language-agnostic BERT sentence
embedding. \emph{ACL} (2022).

\bibitem{bertscore2020} T.~Zhang et al. BERTScore: Evaluating text
generation with BERT. \emph{ICLR} (2020).

\bibitem{cider2015} R.~Vedantam, C.~L.~Zitnick, D.~Parikh. CIDEr:
Consensus-based image description evaluation. \emph{CVPR} (2015).

\bibitem{bleu2002} K.~Papineni, S.~Roukos, T.~Ward, W.-J.~Zhu. BLEU: a
method for automatic evaluation of machine translation. \emph{ACL}
(2002).

\bibitem{xm3600} A.~V.~Thapliyal et al. Crossmodal-3600: A massively
multilingual multimodal evaluation dataset. \emph{EMNLP} (2022).

\end{thebibliography}

\end{document}
